\chapter{The SSH Chain and Two-Band Chiral Fermions}
\label{chap:ssh-chain-two-band-chiral-fermions}

The Su--Schrieffer--Heeger (SSH) chain is not a Jordan--Wigner image of the XY
mother model.  It is nevertheless an essential companion example because it is
the simplest number-conserving two-band topological Hamiltonian.  The Kitaev
chain uses a Nambu particle--hole basis; the SSH chain uses a sublattice basis.
Both reduce topology to the winding of a two-component vector around the origin,
but the Hilbert-space meaning of the two components is different.

\section{SSH Hamiltonian}
\label{sec:ssh-hamiltonian}

Consider a one-dimensional chain with two sites per unit cell, labelled $A$ and
$B$.  Let $a_j$ annihilate a fermion on sublattice $A$ in unit cell $j$, and let
$b_j$ annihilate a fermion on sublattice $B$.  The SSH Hamiltonian is
\begin{equation}
\label{eq:ssh-real-space-hamiltonian}
    H_{\rm SSH}
    =
    \sum_{j}
    \left(
        t_1 a_j^\dagger b_j
        +
        t_2 b_j^\dagger a_{j+1}
        +
        \text{h.c.}
    \right).
\end{equation}
Here $t_1$ is the intracell hopping and $t_2$ is the intercell hopping.  The
model conserves fermion number:
\begin{equation}
\label{eq:ssh-number-conservation}
    [H_{\rm SSH},N]=0,
    \qquad
    N=\sum_j(a_j^\dagger a_j+b_j^\dagger b_j).
\end{equation}
Thus it is a Bloch-band problem rather than a BdG problem.

The two exactly solvable limits are instructive.  If $t_2=0$, the chain is a set
of isolated intracell dimers.  If $t_1=0$, the chain is a set of intercell dimers;
for an open chain, the leftmost $A$ site and rightmost $B$ site are unpaired edge
states.  The topological distinction between these two dimerizations is the
central content of the SSH model.

\section{Bloch form}
\label{sec:ssh-bloch-form}

Use the Fourier convention
\begin{equation}
    a_j=\frac{1}{\sqrt{L}}\sum_k\ee^{\ii kj}a_k,
    \qquad
    b_j=\frac{1}{\sqrt{L}}\sum_k\ee^{\ii kj}b_k.
\end{equation}
In the basis
\begin{equation}
\label{eq:ssh-bloch-spinor}
    \psi_k=
    \begin{pmatrix}
        a_k\\ b_k
    \end{pmatrix},
\end{equation}
the Hamiltonian becomes
\begin{equation}
\label{eq:ssh-momentum-hamiltonian}
    H_{\rm SSH}
    =
    \sum_k\psi_k^\dagger h_{\rm SSH}(k)\psi_k,
\end{equation}
where
\begin{equation}
\label{eq:ssh-bloch-matrix}
    h_{\rm SSH}(k)
    =
    \begin{pmatrix}
        0&q(k)\\
        q(k)^*&0
    \end{pmatrix},
    \qquad
    q(k)=t_1+t_2\ee^{-\ii k}.
\end{equation}
Equivalently,
\begin{equation}
\label{eq:ssh-d-vector}
    h_{\rm SSH}(k)
    =
    d_x(k)\sigma_x+d_y(k)\sigma_y,
\end{equation}
with
\begin{equation}
\label{eq:ssh-d-components}
    d_x(k)=t_1+t_2\cos k,
    \qquad
    d_y(k)=t_2\sin k.
\end{equation}
The energy bands are
\begin{equation}
\label{eq:ssh-band-energies}
    E_\pm(k)=\pm\abs{q(k)}
    =
    \pm\sqrt{t_1^2+t_2^2+2t_1t_2\cos k}.
\end{equation}
The gap closes when
\begin{equation}
\label{eq:ssh-gap-closing}
    q(k)=0.
\end{equation}
For real positive hoppings this occurs at
\begin{equation}
\label{eq:ssh-critical-point}
    t_1=t_2,
    \qquad
    k=\pi.
\end{equation}

\section{Chiral symmetry and winding number}
\label{sec:ssh-chiral-symmetry-winding}

The SSH Bloch Hamiltonian anticommutes with $\sigma_z$:
\begin{equation}
\label{eq:ssh-chiral-symmetry}
    \sigma_z h_{\rm SSH}(k)\sigma_z=-h_{\rm SSH}(k).
\end{equation}
This is sublattice, or chiral, symmetry.  It forces the Hamiltonian to be
off-diagonal in the $(A,B)$ basis and pins the spectrum symmetrically around zero
energy.

When the bulk is gapped, $q(k)$ never vanishes, so the map
\begin{equation}
\label{eq:ssh-q-map}
    k\in S^1_{\rm BZ}
    \longmapsto
    q(k)\in\CC\setminus\{0\}
\end{equation}
has an integer winding number.  With the off-diagonal convention
\cref{eq:ssh-bloch-matrix}, define
\begin{equation}
\label{eq:ssh-winding-number}
    \nu_{\rm SSH}
    =
    \frac{1}{2\pi\ii}
    \int_{-\pi}^{\pi}\dd k\,
    \partial_k\log q(k).
\end{equation}
For $t_1,t_2>0$ and $q(k)=t_1+t_2\ee^{-\ii k}$, the result is
\begin{equation}
\label{eq:ssh-winding-result}
    \nu_{\rm SSH}
    =
    \begin{cases}
        -1, & t_2>t_1,\\
        0, & t_1>t_2.
    \end{cases}
\end{equation}
The sign is an orientation convention coming from the choice $\ee^{-\ii k}$ in
$q(k)$.  Replacing $q(k)$ by $t_1+t_2\ee^{\ii k}$ reverses the sign.  The robust
physical distinction is whether the winding is zero or nonzero.

Geometrically, $q(k)$ traces a circle of radius $\abs{t_2}$ centered at $t_1$ in
the complex plane.  The nontrivial phase is precisely the case in which this
circle encloses the origin:
\begin{equation}
\label{eq:ssh-topological-condition}
    \abs{t_2}>\abs{t_1}.
\end{equation}

\section{Zak phase}
\label{sec:ssh-zak-phase}

Let
\begin{equation}
    q(k)=\abs{q(k)}\ee^{\ii\varphi(k)}.
\end{equation}
A normalized occupied-band eigenstate may be chosen locally as
\begin{equation}
\label{eq:ssh-occupied-eigenstate}
    \ket{u_-(k)}
    =
    \frac{1}{\sqrt{2}}
    \begin{pmatrix}
        -\ee^{\ii\varphi(k)}\\ 1
    \end{pmatrix}.
\end{equation}
Its Berry connection is
\begin{equation}
\label{eq:ssh-berry-connection}
    \mathcal{A}(k)
    =
    \ii\braket{u_-(k)}{\partial_k u_-(k)}
    =
    -\frac12\partial_k\varphi(k),
\end{equation}
up to gauge conventions.  The Zak phase is therefore
\begin{equation}
\label{eq:ssh-zak-phase}
    \Phi_{\rm Zak}
    =
    \int_{-\pi}^{\pi}\mathcal{A}(k)\,\dd k
    =
    -\pi\nu_{\rm SSH}
    \quad
    \text{mod }2\pi.
\end{equation}
Thus the nontrivial SSH phase has Zak phase $\pi$ modulo $2\pi$, provided the
unit-cell convention is fixed.  Changing the unit cell shifts the Zak phase by an
integer multiple of $2\pi$ or changes the representative convention, but it does
not remove the protected boundary-state distinction for a specified physical
termination.

\section{Open chain and edge modes}
\label{sec:ssh-edge-modes}

For an open chain with $L$ unit cells, the zero-energy eigenvalue equation at
energy $E=0$ decouples between sublattices.  A left edge mode supported on the
$A$ sublattice has the form
\begin{equation}
\label{eq:ssh-left-edge-ansatz}
    \ket{\psi_L}
    =
    \sum_{j=1}^{L}\alpha_j a_j^\dagger\ket{0}.
\end{equation}
The Schrödinger equation gives the recursion
\begin{equation}
\label{eq:ssh-edge-recursion}
    t_1\alpha_j+t_2\alpha_{j+1}=0,
\end{equation}
so that
\begin{equation}
\label{eq:ssh-edge-wavefunction}
    \alpha_j
    =
    \left(-\frac{t_1}{t_2}\right)^{j-1}\alpha_1.
\end{equation}
This mode is normalizable at the left edge in the thermodynamic limit if and only
if
\begin{equation}
\label{eq:ssh-edge-normalizability}
    \abs{t_1}<\abs{t_2}.
\end{equation}
The right edge mode is similarly localized on the $B$ sublattice.  Thus the
bulk winding condition and the edge-state condition agree.

At the exactly dimerized point $t_1=0$, the edge modes are simply the unpaired
operators $a_1^\dagger$ and $b_L^\dagger$.  This is the SSH analogue of the
unpaired Majorana operators in the topological limit of the Kitaev chain.

\section{Relation and distinction from the Kitaev chain}
\label{sec:ssh-relation-distinction-kitaev}

Both the SSH and Kitaev chains are two-band topological models, but the meanings
of the two-component spinors are different:
\begin{equation}
\label{eq:ssh-kitaev-spinor-comparison}
    \psi_k^{\rm SSH}
    =
    \begin{pmatrix}
        a_k\\ b_k
    \end{pmatrix},
    \qquad
    \Psi_k^{\rm Kitaev}
    =
    \begin{pmatrix}
        c_k\\ c_{-k}^\dagger
    \end{pmatrix}.
\end{equation}
The SSH Pauli matrices act in sublattice space.  The Kitaev Pauli matrices act in
Nambu particle--hole space.  Consequently:
\begin{itemize}[leftmargin=2.2em]
    \item the SSH chain conserves fermion number, while the Kitaev chain
    conserves only fermion parity;
    \item the SSH spectral symmetry comes from sublattice chiral symmetry,
    while the Kitaev BdG spectral symmetry comes from particle--hole redundancy;
    \item SSH edge states are ordinary complex-fermion boundary orbitals, while
    Kitaev edge states are Majorana zero modes;
    \item the SSH winding is naturally a chiral-band winding, while the Kitaev
    winding is a BdG/Nambu winding, reducible to a class-D $\ZZ_2$ invariant when
    the real-pairing chiral symmetry is not imposed.
\end{itemize}

This comparison is useful because the same mathematical form
\begin{equation}
\label{eq:ssh-kitaev-two-band-common-form}
    H(k)=d_x(k)\sigma_x+d_y(k)\sigma_y+d_z(k)\sigma_z
\end{equation}
can describe physically distinct objects.  A two-band matrix by itself does not
say whether the Pauli matrices refer to sublattice, spin, orbital, or Nambu
space.  The symmetry interpretation must always be specified.

\section{SSH as a two-band chiral mother example}
\label{sec:ssh-two-band-chiral-mother}

The SSH chain represents the simplest member of a broader family of one-dimensional chiral two-band Hamiltonians:
\begin{equation}
\label{eq:ssh-general-chiral-two-band}
    h(k)=
    \begin{pmatrix}
        0&q(k)\\ q(k)^*&0
    \end{pmatrix}.
\end{equation}
Every such gapped Hamiltonian has a winding number
\begin{equation}
\label{eq:ssh-general-winding}
    \nu
    =
    \frac{1}{2\pi\ii}\int_{\rm BZ}\dd k\,
    \partial_k\log q(k).
\end{equation}
Longer-range hoppings simply replace $q(k)=t_1+t_2\ee^{-\ii k}$ by a finite
Laurent polynomial in $\ee^{\ii k}$.  The topology is still controlled by whether
and how this complex loop winds around the origin.  The SSH model is the minimal
case in which this winding can be nonzero.

\section{Summary}
\label{sec:ssh-summary}

The SSH chain is a number-conserving two-band model with Bloch Hamiltonian
\begin{equation}
    h_{\rm SSH}(k)=
    \begin{pmatrix}
        0&t_1+t_2\ee^{-\ii k}\\
        t_1+t_2\ee^{\ii k}&0
    \end{pmatrix}.
\end{equation}
It has chiral symmetry and an integer winding number.  For real positive
hoppings, the nontrivial phase occurs when $t_2>t_1$ and supports boundary zero
modes on an open chain.  It parallels the Kitaev chain at the level of two-band
topology, but differs physically because its two-component structure is
sublattice space rather than Nambu space.
